\section*{Chapitre \ref{ch:ll:queues-ring-buffers-and-shared-memory} --- Files, tampons circulaires et mémoire partagée}

\begin{solution}{exo:ll:queues-ring-buffers-and-shared-memory:1}
$70\,001 - 69\,500 = 501$ messages attendent. L'écriture suivante va dans l'emplacement $70\,001 \bmod 1\,024 = 369$. Le tampon est plein à 1\,024 en attente,
donc $1\,024 - 501 = 523$ de plus.
\end{solution}

\begin{solution}{exo:ll:queues-ring-buffers-and-shared-memory:2}
$2^{64}/(2.5 \times 10^8)$ secondes représentent environ 2\,338 ans~: jamais. Une position sur 32 bits reboucle en 17.2 secondes, ce qui explique pourquoi les positions sont
sur 64 bits et seul l'indice d'emplacement est pris modulo la taille.
\end{solution}

\begin{solution}{exo:ll:queues-ring-buffers-and-shared-memory:3}
$0.8^n < 10^{-6}$ pour $n \ge 62$~; la puissance de deux suivante est 64 emplacements.
\end{solution}

\begin{solution}{exo:ll:queues-ring-buffers-and-shared-memory:4}
Lire la position de l'autre côté revient à lire une ligne que l'autre cœur écrit à chaque message, si bien qu'elle passe d'un cœur à
l'autre à chaque fois. Avec les copies, chaque côté ne lit la ligne de l'autre que lorsque le tampon paraît plein ou vide~; sans elles, chaque message ajoute une
ou deux transferts de ligne, des dizaines de nanosecondes sous contention.
\end{solution}

\begin{solution}{exo:ll:queues-ring-buffers-and-shared-memory:5}
Le retard croît de $3\,000\,000 - 400\,000 = 2.6$ millions par seconde et atteint 131\,072 après 50.4 millisecondes~: le moniteur est
dépassé environ 30 millisecondes avant la fin de la rafale.
\end{solution}

\begin{solution}{exo:ll:queues-ring-buffers-and-shared-memory:6}
Dans le modèle mémoire de C++, une lecture non atomique qui entre en course avec une écriture est une \omterm{def:ll:rust-for-low-latency:race}{course aux données}, c'est-à-dire un \omterm{def:ll:cpp-for-latency-iii-the-compiler:ub}{comportement indéfini} quel que soit l'usage fait ensuite
de la valeur~; le compilateur peut supposer que cela ne se produit pas. Les lectures et écritures atomiques relâchées des mots rendent la course définie,
et la vérification de séquence rend le résultat cohérent.
\end{solution}

\begin{solution}{exo:ll:queues-ring-buffers-and-shared-memory:7}
Les deux positions se retrouvent alors sur une même ligne, que le producteur écrit à chaque message et le consommateur à chaque message~: \omterm{def:ll:memory-hierarchy-and-caches:false}{faux partage} (false sharing)
(chapitre~3). Le débit devrait chuter, car chaque écriture invalide la copie de la ligne détenue par l'autre cœur, tandis que le temps d'aller-retour change
moins, car dans un échange en aller-retour les lignes migrent de toute façon~; la mesure tranche sur l'ampleur de l'écart.
\end{solution}

\begin{solution}{exo:ll:queues-ring-buffers-and-shared-memory:8}
Un pointeur est une adresse dans l'espace d'adressage d'un processus~; il ne signifie rien, ou autre chose, dans l'autre. Une
\texttt{std::string} contient un tel pointeur vers une mémoire qui ne se trouve pas du tout dans la région. Les dispositions partagées ne contiennent que des décalages et des champs de taille fixe
(un tableau de caractères de taille fixe pour le nom), accompagnés d'un numéro de version.
\end{solution}

\begin{solution}{pb:ll:queues-ring-buffers-and-shared-memory:1}
\begin{enumerate}
\item $65\,536 \times 64 = 4$~Mio auxquels s'ajoute l'en-tête~: le double d'un cache de deuxième niveau de 2~Mio~; la partie chaude du tampon est constituée de ses emplacements récents.
\item 7.5~\%, 12~\% et 20~\%.
\item Le moniteur, seul consommateur plus lent que 2 millions par seconde.
\item La bourse n'attendra pas~: un producteur bloqué sur son lecteur le plus lent laisserait plutôt tomber des paquets au niveau de la socket.
\item $65\,536/(2\,000\,000 - 1\,500\,000) = 131$ millisecondes.
\item Oui, après 131 millisecondes d'une rafale de 500 millisecondes.
\item Au moins $0.5 \times 500\,000 = 250\,000$ emplacements~: 262\,144 en puissance de deux.
\item 16~Mio.
\item 5\,000 lignes de 64 octets, soit 320\,000 octets~; à 2 millions de mises à jour par seconde contre 1.5 million de lectures, il en saute un quart,
ne voyant que la plus récente.
\item Il signale le trou, saute à une position récente et reconstruit sa vue (à partir de son propre instantané ou de celui du constructeur de carnet).
\item Perdre des données de marché est récupérable (instantanés, \omterm{def:ll:queues-ring-buffers-and-shared-memory:bp}{conflation}), alors qu'attendre ne l'est pas~; perdre un ordre ou le réordonner ne l'est ni l'un ni l'autre, si bien que la
voie des ordres attend ou refuse.
\item $\P(N \ge n) = 0.12^n$, négligeable pour tout $n$ raisonnable~; optimiste, car les arrivées surviennent par rafales, bien loin d'un processus de Poisson.
\item \textbf{Résultat nommé.} Le moniteur est dépassé 131 millisecondes après le début d'une rafale de 2 millions par seconde sur un tampon de 65\,536 emplacements~; en lisant
plutôt à l'aide de \omterm{def:ll:queues-ring-buffers-and-shared-memory:seqlock}{verrous à séquence} (seqlock) par instrument, il saute 25~\% des mises à jour et n'est jamais dépassé.
\item Environ \qty{200}{\nano\second} par aller-retour sur cet ordinateur portable~: un consommateur ajoute un curseur, non une file, et son coût est celui de sa propre lecture.
\item Parce que deux binaires compilés à des moments différents, ou dans des langages différents, la lisent~: un changement doit être détecté, non mal interprété.
\item Des centaines de fois la latence du tampon par saut, avec une queue de distribution en millisecondes, comme l'a montré la comparaison publiée.
\item Tous deux sur le nœud de la carte réseau, sur des cœurs qui se transmettent les lignes le plus vite (chapitre~4).
\item Enregistrer le taux d'arrivée par milliseconde à partir de la capture au niveau du fil (chapitre~23) et conserver la distribution des longueurs de rafales
et des pics.
\item Invalider l'état touché et attendre, ou demander, une resynchronisation, comme le fait le constructeur de carnet en cas de trou (chapitre~19).
\item Les rafales qu'il doit absorber.
\end{enumerate}
\end{solution}

\begin{solution}{iq:ll:queues-ring-buffers-and-shared-memory:1}
Un tableau de taille puissance de deux, une position de producteur et une position de consommateur sur des \omterm{def:ll:memory-hierarchy-and-caches:line}{lignes de cache} distinctes, chacune écrite par un seul côté. Le producteur
écrit l'emplacement, puis enregistre sa position avec release~; le consommateur la charge avec acquire, lit, puis enregistre la sienne avec release,
que le producteur charge avec acquire avant de réutiliser des emplacements. On met en cache la position de l'autre côté.
\iqlookfor{paires release/acquire, lignes distinctes, positions mises en cache.}
\end{solution}

\begin{solution}{iq:ll:queues-ring-buffers-and-shared-memory:2}
Pas de verrou pour lequel rivaliser ni de sommeil~: le consommateur tourne sur une position présente dans son cache et voit l'écriture du producteur sous forme de
transfert de \omterm{def:ll:memory-hierarchy-and-caches:line}{ligne de cache}, une centaine de nanosecondes environ, là où une variable de condition endort le thread dans le noyau et le réveille par
un appel système et un changement de contexte, soit des dizaines de microsecondes.
\iqlookfor{intervention du noyau face au coût de \omterm{def:ll:memory-hierarchy-and-caches:coherence}{cohérence de cache}.}
\end{solution}

\begin{solution}{iq:ll:queues-ring-buffers-and-shared-memory:3}
Un compteur mis à une valeur impaire pendant une écriture et à une valeur paire après~; les lecteurs copient puis recommencent si le compteur a changé ou était impair. Convient aux petites valeurs
avec un seul écrivain et de nombreux lecteurs qui ne veulent que la plus récente (sommet du carnet). Pièges~: la copie doit utiliser des lectures atomiques en C++ pour éviter une \omterm{def:ll:rust-for-low-latency:race}{course
aux données}~; les lecteurs peuvent être affamés sous des écritures constantes~; l'écrivain doit être unique.
\iqlookfor{le protocole de nouvelle tentative et la \omterm{def:ll:rust-for-low-latency:race}{course aux données} au niveau du langage.}
\end{solution}

\begin{solution}{iq:ll:queues-ring-buffers-and-shared-memory:4}
Un seul tampon de diffusion écrit par le feed handler, chaque consommateur à sa propre position, le producteur n'attendant jamais~; les consommateurs qui
ne suivent pas reçoivent des vues conflatées (dernière valeur par instrument) ou détectent les dépassements et se resynchronisent à partir d'un instantané~; les tailles sont fixées
d'après les rafales enregistrées.
\iqlookfor{une politique pour le consommateur le plus lent.}
\end{solution}

\begin{solution}{iq:ll:queues-ring-buffers-and-shared-memory:5}
Les pointeurs et tout ce qui en contient (chaînes, conteneurs, tables virtuelles), les champs dont la taille ou l'alignement dépend du
compilateur, les dispositions sans version~; les verrous qui ne peuvent pas être récupérés si le détenteur meurt.
\iqlookfor{indépendance vis-à-vis de l'espace d'adressage et versionnement.}
\end{solution}

\begin{solution}{iq:ll:queues-ring-buffers-and-shared-memory:6}
Partir de la queue M/M/1 $\P(N \ge n) = \rho^n$ pour une probabilité de dépassement visée~; vérifier ensuite sur les rafales enregistrées, qui sont loin
d'un processus de Poisson et durent quelques secondes~; dimensionner pour la pire rafale augmentée d'une marge, et concevoir ce qui se passe en cas de dépassement. La théorie des files d'attente donne
l'ordre de grandeur et l'effet du taux d'utilisation~; elle ne voit ni la corrélation ni les rafales.
\iqlookfor{la théorie pour la forme, les données pour la taille.}
\end{solution}
